\documentclass[11pt]{article}
\usepackage[margin=1in]{geometry}
\usepackage{amsmath,amssymb,amsthm,mathtools}
\usepackage{graphicx}
\usepackage[hidelinks]{hyperref}
\emergencystretch=3em

\newtheorem{theorem}{Theorem}
\newtheorem{lemma}{Lemma}
\newtheorem{remark}{Remark}
\newcommand{\D}{\mathbb D}
\newcommand{\T}{\mathbb T}
\newcommand{\dm}{\,dm}
\newcommand{\Hp}{H^p}
\newcommand{\Bop}{\mathbf B}

\title{Explicit Two-Sided Bounds for the Backward Shift on $H^6$}
\author{Run fa\_banach\_001, lane 10}
\date{August 11, 2026}

\begin{document}
\maketitle

\section*{Source Question}

Karlovych and Shargorodsky ask in Open Problem 5.4 of
arXiv:2308.04072 to determine the exact norm of the backward shift
\[
  (\Bop f)(z)=\frac{f(z)-f(0)}{z}
\]
on $H^p$ for $1\le p<\infty$ \cite{KS}.  The source passage is reproduced in
Figure~\ref{fig:open}.

\begin{figure}[t]
\centering
\includegraphics[width=.94\textwidth]{figures/open_problem_crop.png}
\caption{Open Problem 5.4 in the source PDF, page 15.}
\label{fig:open}
\end{figure}

The problem has subsequently been solved at $p=1$ by Brevig and Seip
\cite{BS}, and at $p=4$ by Bampouras and Llinares \cite{BL}; of course the
values at $p=2$ and $p=\infty$ were already known.  This packet gives a
rigorous explicit bracket at the next even exponent $p=6$.  It does not claim
the exact $H^6$ norm.

\section*{Definitions}

Normalized arclength measure on $\T$ is denoted by $m$.  For $1\le p<\infty$,
\[
  \|f\|_{H^p}^p=\int_{\T}|f(\zeta)|^p\dm(\zeta).
\]
Write $\varphi=(1+\sqrt5)/2$.  We use the known sharp values
\[
 \|\Bop\|_{H^1}=\frac2{\sqrt3},\qquad
 \|\Bop\|_{H^2}=1,\qquad
 \|\Bop\|_{H^4}=\varphi^{1/4},\qquad
 \|\Bop\|_{H^\infty}=2.
\]

\section*{Proof Intuition}

For the lower bound, test the operator on an affine function of a disk
automorphism $I$.  Because $I$ is inner, every mixed moment
$\int I^j\overline I^k$ reduces to a power of $I(0)$.  At $p=6$ this makes
both norms exact degree-six polynomials in two rational parameters.  A
near-optimal rational choice gives the stated exact lower bound.

For the upper bound, the Hardy spaces form the usual complex interpolation
scale.  Interpolating the new sharp $H^4$ norm with the sharp $H^\infty$ norm
gives a global estimate for $4\le p\le\infty$, and hence the $H^6$ upper
bound.  The same argument between the other known endpoints gives the
piecewise estimate recorded below.

\section*{The Inner-Function Calculation}

\begin{lemma}\label{lem:moments}
For $0\le a<1$, let
\[
  I_a(z)=\frac{a+z}{1+az}.
\]
For real $c$, define $D(a,c)=\|I_a-c\|_{H^6}^6$.  Then
\begin{align*}
D(a,c)={}&1+9c^2+9c^4+c^6-6ac-18ac^3-6ac^5\\
&\quad+6a^2c^2+6a^2c^4-2a^3c^3,
\end{align*}
and
\[
  D(a,a)=(1-a^2)(a^4+4a^2+1).
\]
\end{lemma}

\begin{proof}
The function $I_a$ is inner and $I_a(0)=a$.  Hence for integers
$0\le j,k\le3$,
\[
 \int_{\T}I_a^j\overline{I_a}^{,k}\dm=a^{|j-k|}.
\]
Indeed, if $j\ge k$, unimodularity on $\T$ reduces the integrand to
$I_a^{j-k}$, whose mean is $I_a(0)^{j-k}$; the other case follows by
conjugation.  Expanding
\[
  D(a,c)=\|(I_a-c)^3\|_{H^2}^2
\]
and inserting these moments gives the displayed polynomial.  Substitution
$c=a$ and factorization gives the second identity.
\end{proof}

\section*{Main Partial Result}

\begin{theorem}[Explicit $H^6$ bracket]\label{thm:h6}
The backward shift on $H^6$ satisfies
\[
 \left(\frac{6229716996096}{1533633162625}\right)^{1/6}
 \le \|\Bop\|_{H^6}
 \le 2^{1/3}\varphi^{1/6}.
\]
Numerically,
\[
  1.2631583709 \le \|\Bop\|_{H^6}\le1.3651318687.
\]
\end{theorem}

\begin{proof}
For the lower bound take
\[
 a=\frac{32}{35},\qquad c=\frac{11}{28},\qquad f=I_a-c.
\]
Since $f(0)=a-c$ and $|z|=1$ on $\T$,
\[
  \|\Bop f\|_{H^6}^6
  =\|I_a-a\|_{H^6}^6=D(a,a).
\]
Lemma~\ref{lem:moments} and exact rational arithmetic give
\[
 \frac{\|\Bop f\|_{H^6}^6}{\|f\|_{H^6}^6}
 =\frac{D(32/35,32/35)}{D(32/35,11/28)}
 =\frac{6229716996096}{1533633162625}.
\]
Taking sixth roots proves the lower estimate.

For the upper bound, use the standard complex interpolation identity for
Hardy spaces \cite[Chapter 7]{Duren}.  Interpolating $\Bop:H^4\to H^4$,
whose sharp norm is $\varphi^{1/4}$ by \cite{BL}, with
$\Bop:H^\infty\to H^\infty$, whose norm is $2$, yields
\[
 \|\Bop\|_{H^p}\le
 (\varphi^{1/4})^{4/p}2^{1-4/p}
 =\varphi^{1/p}2^{1-4/p},\qquad 4\le p\le\infty.
\]
Putting $p=6$ gives the claimed upper estimate.
\end{proof}

\begin{theorem}[Interpolation bounds for all finite exponents]
For $1\le p<\infty$,
\[
\|\Bop\|_{H^p}\le
\begin{cases}
 (2/\sqrt3)^{2/p-1},&1\le p\le2,\\[2mm]
 \varphi^{1/2-1/p},&2\le p\le4,\\[2mm]
 \varphi^{1/p}2^{1-4/p},&4\le p<\infty.
\end{cases}
\]
\end{theorem}

\begin{proof}
Interpolate respectively between the sharp endpoint pairs
$(H^1,H^2)$, $(H^2,H^4)$, and $(H^4,H^\infty)$.  Solving
$1/p=(1-\theta)/p_0+\theta/p_1$ and applying the operator interpolation
inequality gives exactly the three displayed exponents.
\end{proof}

\section*{Verification and Upgrade Attempts}

The script \texttt{code/verify\_h6\_bound.py} independently expands the mixed
moments, verifies the factorization of $D(a,a)$, checks the exact rational
quotient, and prints the numerical endpoints.

Several routes toward the exact $H^6$ value were tested before stopping:
\begin{enumerate}
\item A current-literature check found the $H^1$ and May 2026 $H^4$ solutions,
but no $H^6$ solution.
\item Optimizing the affine-inner family numerically gives
$a\approx0.9146566921$, $c\approx0.3930991838$, and sixth-power ratio
$4.0620742602$; the rational pair in Theorem~\ref{thm:h6} is within
$10^{-5}$ at the sixth-power level.
\item A degree-six harmonic moment-majorant was sought with equality on the
candidate circle.  The resulting polynomial is divisible by the circle
equation only once, not twice, so it cannot be globally nonnegative.  Thus the
natural scalar/SOS certificate cannot prove exactness.
\item Direct Euler--Lagrange expansion at $p=6$ introduces three independent
mixed correlations; the single sharp Fourier-symmetry lemma that closes the
$H^4$ proof does not control them simultaneously.
\item Optimized finite Taylor polynomials of degrees up to six remained below
the affine-inner value and therefore supplied evidence but no upper bound.
\item Powers $(I-c)^k$, $2\le k\le6$, were numerically optimized and gave
strictly smaller ratios.
\item General real polynomials of degree up to four in one inner function did
not improve the affine family in bounded numerical searches.
\item Interpolation gives the rigorous upper estimate above, but its equality
conditions are incompatible with the known $H^4$ extremals and do not identify
an $H^6$ extremizer.
\end{enumerate}
These obstructions leave a genuine gap of about $0.102$ in operator norm, so
the exact $H^6$ problem remains open in this packet.

\section*{Scope and Novelty Check}

This is a partial answer to Open Problem 5.4, not a determination of the exact
norm.  The exact cases currently identified are $p=1,2,4,\infty$.  Cheap run
indexes were searched for arXiv:2308.04072, backward shift, $H^p$, and Hardy
space norm.  Web/arXiv searches on August 11, 2026 used the exact phrases
``norm of the backward shift operator on $H^p$'', ``backward shift on $H^4$'',
and ``backward shift $H^6$ exact norm''.  They found \cite{BS,BL} but no exact
$H^6$ result.

The May 2026 paper \cite{BL} announces a forthcoming preprint by Jim\'enez and
Vukoti\'c on general norm estimates.  That preprint was not findable by exact
title at the time of this check, so overlap with the explicit lower-bound
family remains possible.  Novelty confidence is therefore moderate; proof
validity does not depend on novelty.

\section*{Human-Review Recommendation}

The core proof is elementary.  Review should check the mixed-moment expansion
in Lemma~\ref{lem:moments}, the exact rational substitution, and use of the
standard Hardy-space interpolation scale including the $H^\infty$ endpoint.

\begin{thebibliography}{9}
\bibitem{KS} O. Karlovych and E. Shargorodsky,
\emph{Bounded compact and dual compact approximation properties of Hardy
spaces: new results and open problems}, arXiv:2308.04072 (2023).
\bibitem{BS} O. F. Brevig and K. Seip,
\emph{The norm of the backward shift on $H^1$ is $2/\sqrt3$},
Pure Appl. Funct. Anal. 9 (2024), 991--994; arXiv:2309.11360.
\bibitem{BL} K. Bampouras and A. Llinares,
\emph{The norm of the backward shift on $H^4$ is $\sqrt[4]{\varphi}$},
arXiv:2605.10469 (2026).
\bibitem{Duren} P. L. Duren, \emph{Theory of $H^p$ Spaces}, Academic Press,
1970.
\end{thebibliography}

\end{document}
